%%%

%

%\documentclass[11pt,a4paper,twoside]{article}

\usepackage[utf8]{inputenc}

\usepackage[T1,T2A]{fontenc}

\usepackage[english.ancient,russian.ancient]{babel}

\usepackage{graphicx}

\usepackage{array}

\usepackage[doi]{samgtu-bib}

\usepackage{samgtu}

\usepackage{samgtu-name}

\usepackage{mathtools}

\mathtoolsset{showonlyrefs}

\usepackage{desclist}

\usepackage{euscript}

\usepackage{mathrsfs} 

\usepackage{multicol}

\usepackage{mathrsfs}

\usepackage{cite}

\usepackage{cmap}

\usepackage{qrcode}

\allowdisplaybreaks[4]

\usepackage[nodayofweek]{datetime}

\usepackage[thinlines]{easybmat}



\renewcommand{\JournalYear}{2018}

\renewcommand{\JournalVolume}{22}

\renewcommand{\JournalNumber}{4}



\newdate{JournalDateSign}{29}{12}{2018}% Дата подписания Вестника в печать

\renewcommand{\JournalDateSign}{\displaydate{JournalDateSign}}

\newdate{JournalDatePrint}{16}{01}{2019}% Дата выхода Вестника из печати

\renewcommand{\JournalDatePrint}{\displaydate{JournalDatePrint}}

%\renewcommand{\JournalUPL}{16.56} % Усл. печ. л.

%\renewcommand{\JournalUIL}{16.53} % Уч.-изд. л.

\renewcommand{\JournalUPL}{15.24} % Усл. печ. л.

\renewcommand{\JournalUIL}{15.21} % Уч.-изд. л.

\renewcommand{\JournalRegNumber}{220/18} % Регистрационный номер

\renewcommand{\JournalOrderNumber}{2} % Номер заказа



%\renewcommand{\JournalRMonth}{Март} % Месяц выхода Вестника

%\renewcommand{\JournalEMonth}{March} % Месяц выхода Вестника

%\renewcommand{\JournalRMonth}{Июнь} % Месяц выхода Вестника

%\renewcommand{\JournalEMonth}{June} % Месяц выхода Вестника

%\renewcommand{\JournalRMonth}{Сентябрь} % Месяц выхода Вестника

%\renewcommand{\JournalEMonth}{September} % Месяц выхода Вестника

\renewcommand{\JournalRMonth}{Декабрь} % Месяц выхода Вестника

\renewcommand{\JournalEMonth}{December} % Месяц выхода Вестника



\newcommand{\totop}[1]{\raisebox{6pt}[1pt][2pt]{#1}}

\newcommand{\tottop}[1]{\raisebox{12pt}[1pt][2pt]{#1}} 

\newcommand{\totttop}[1]{\raisebox{18pt}[1pt][2pt]{#1}} 

\newcommand{\tottttop}[1]{\raisebox{24pt}[1pt][2pt]{#1}} 

\newcommand{\totttttop}[1]{\raisebox{30pt}[1pt][2pt]{#1}} 

\newcommand{\tobot}[1]{\raisebox{-6pt}[1pt][2pt]{#1}} 

\newcommand{\tobbot}[1]{\raisebox{-12pt}[1pt][2pt]{#1}} 

\newcommand{\tobbbot}[1]{\raisebox{-18pt}[1pt][2pt]{#1}}



\graphicspath{{./00PIC/}}

%

%\usepackage{samgtu-name}

%\usepackage{wrapfig}

%

%

\vsgtu{1605}

%

%

\Personalia

%\pdfcrop % обрезка pdf для размещения на math-net.ru

%\draft % На корректуре

%

%

%

\FirstPage{7}

\LastPage{14}

%

%

%\makeatletter

%\newcommand{\PersonRadRef}{\@langrustrue\def\refname{{\bf\small Избранные  труды Е.~В.~Радкевича за последние 5 лет}\footnote{Список подготовлен М.~Н.~Саушкиным (Самара, СамГТУ).}}}

%\newcommand{\NoRef}{\@langrustrue\def\refname{{\bf\footnotesize Цикл работ:}}}

%%

%%\newcommand{\ManzhirovRef}{\@langrustrue\def\refname{{\bf\small Статьи о проф. А.~В.~Манжирове}}}

%%\newcommand{\ManzhirovRefEng}{\@langrustrue\def\refname{{\bf\small Publications  about Prof. A.~V.~Manzhirov in the Press}}}

%%

%%\makeatother







%\begin{document}

$\;$



\hfill \qrcode[hyperlink,height=.8in]{http://doi.org/10.14498/vsgtu1605}

\vspace{-.8in}







\rutitle{К 75-летию профессора Евгения Владимировича Радкевича }

\rutitleshort{К 75-летию  профессора Евгения Владимировича Радкевича }

\rutitlehack{К 75-летию профессора\\ Евгения Владимировича Радкевича }

\entitle{To the 75$^{\rm th}$ Anniversary of Professor Evgeniy Vladimirovich Radkevich}

\entitlehack{To the 75$^{\bf th}$ Anniversary of\\ Professor Evgeniy Vladimirovich Radkevich}



\ruauthor{А.~А.~Андреев, В.~П.~Радченко, Е.~А.~Козлова}{А\,н\,д\,р\,е\,е\,в~А.~А., Р\,а\,д\,ч\,е\,н\,к\,о~В.~П., К\,о\,з\,л\,о\,в\,а~Е.~А.}

\enauthor{A.~A.~Andreyev, V.~P.~Padchenko, E.~A.~Kozlova}{A\,n\,d\,r\,e\,y\,e\,v~A.~A., R\,a\,d\,c\,h\,e\,n\,k\,o~V.~P., K\,o\,z\,l\,o\,v\,a~E.~A.}





\ruaffil{\rusamgtu.}

\enaffil{\ensamgtu.}



\selectlanguage{russian}%

\phantomsection%

\addtocontent{{\it\RuAuthorInContents}\hspace{2.5mm}}{``\RuTitleInContents''}

\addtoengcontent{{\it\EnAuthorInContents}\hspace{2.5mm}}{``\selectlanguage{english}\EnTitleInContents''}

\normalsize 

\normalfont

\chead[\fancyplain{}{\selectlanguage{russian}\theruauthor}]{\fancyplain{}{\selectlanguage{russian} \therutitle}}

\rhead[]{\fancyplain{\RuJournal}{}}

\lhead[\fancyplain{\RuJournal}{}]{}

\cfoot[]{}

\lfoot[\thepage]{}

\rfoot[]{\thepage}

\renewcommand{\plainheadrulewidth}{0.7pt}

\renewcommand{\headrulewidth}{0.7pt}

\pagestyle{fancyplain}

\thispagestyle{plain}

\clear

\hskip-7mm

\begin{minipage}[t]{125mm}

\vskip.5mm

\noindent\theudc

\vskip2mm

\RuTitleInTitlepage

\end{minipage}

\par

\theruauthorsdetails

\begin{minipage}[t]{125mm}

\vskip4mm

{\renewcommand{\bfdefault}{bx}\noindent\small\textbf{\emph{\RuAuthorInTitlepage}}\renewcommand{\bfdefault}{b}}  

\vskip2mm

{\scriptsize \ruaffilname}

\end{minipage}

\vskip2mm





\footnotetext{

\vspace{-1em}\\

\noindent{\normalsize\bf\ruArticleType}\smallskip\par\noindent

\OAlogo~\ccby\  Контент публикуется на условиях  лицензии  \href{https://creativecommons.org/licenses/by/4.0/deed.ru}{Creative Commons Attribution 4.0 International} (\url{https://creativecommons.org/licenses/by/4.0/deed.ru})

\vspace{.3em}\\ 

\RuCitation}





\begin{wrapfigure}{o}{0.46\textwidth} 

\vspace{-3mm}

\centering

\includegraphics[width=0.4\textwidth]{radkevich75.jpg}



\vspace{-8mm}



\end{wrapfigure} 



\footnotesize



\smallskip





26 января 2018~г. исполнилось 75 лет доктору физико-математических наук, профессору Евгению Владимировичу Радкевичу.



В настоящей  заметке приводятся сведения из научной биографии этого крупного ученого, который хорошо известен в~России и~за рубежом. Обсуждаются его работы, которые внесли существенный вклад во многие области науки и техники.





\newdate{radkeviche}{29}{03}{2018}

\renewcommand{\JournalDataOnline}{\displaydate{radkeviche}}



\smallskip



\par\noindent

\theonlinerudate \par\noindent

\rule{68mm}{.7pt}\\



\normalsize



Евгений Владимирович Радкевич родился 26 января 1943 года в городе Челябинске. 

В~1965 году он окончил \ruaffid{595}{механико-математический факультет МГУ}, затем поступил в~аспирантуру под руководством д.ф.-м.н. \rupid{8582}{О.~А.~Олейник}. 

В~1969 году им была защищена кандидатская диссертация на тему <<Гипоэллиптические операторы с переменными коэффициентами>>. 

С~1968 года работал в отделе математических методов механики \ruaffid{755}{Института проблем механики АН СССР} (ныне ИПМех РАН). 

С~1974 года перешел на преподавательскую работу в~\ruaffid{8145}{Московский институт электроники и~математики} (ныне одно из подразделений \ruaffid{7333}{НИУ ВШЭ}), позже "--- в~\ruaffid{10626}{Московский институт радиотехники, электроники и автоматики} (ныне Московский технологический университет). 

В~1988 году Евгений Владимирович защитил докторскую диссертацию <<Эллиптико"=параболические краевые задачи со свободной границей>>, а~в~1992 году ему было присвоено ученое звание профессора. 

В~настоящее время работает профессором кафедры дифференциальных уравнений механико"=математического факультета МГУ.



Е.~В.~Радкевич "--- крупный специалист в качественной теории эволюционных краевых задач. Спектр научных интересов Е.~В.~Радкевича включает в себя также задачи со свободной границей, асимптотические методы, нелинейные уравнения, проблемы фазовых переходов, вопросы существования и~свойств аттракторов для решений дифференциальных уравнений и многое другое. Работы Е.~В.~Радкевича внесли существенный вклад во многие области науки и техники.



В~области уравнений с~частными производными им получены достаточные условия гипоэллиптичности операторов с~кратными характеристиками, необходимые для операторов второго порядка. 

Совместно с~О.~А.~Олейник получены результаты в~теории уравнений второго порядка с~неотрицательной характеристической формой. 

Для предельных задач математических моделей фазовых переходов разработан метод модельных операторов, позволивший получить теоремы существования классических решений для некоторого класса задач со свободной границей. 

Обобщен метод операторных пучков (уравнений Карлемана) для исследования проблемы динамического угла, получены условия существования классического решения проблемы динамического угла для модифицированной задачи Стефана и~задачи Веригина"--~Маскета. 

Построены и~обоснованы асимптотические решения системы фазового поля, неизотермической системы Кана"--~Хилларда и~модели внутренних слоев как регуляризации классической задачи Веригина"--~Маскета. 

Построена и~обоснована асимптотика взаимодействия и~аттракции нелинейных параболических волн в моделях теории фазовых переходов (система фазового поля, система Кана"--~Хилларда, системы Мимури). 

Описаны предельные задачи, отвечающие решениям типа волн Ван дер Вальса, солитона, т.н. волновым поездам (системы внутренних слоев), доказана их классическая разрешимость. 

Получено волновое обоснование <<mushy region>> как теплового пульсара в~системе фазового поля. 





Е.~В.~Радкевич является автором более 100 научных работ и трех монографий. 

Монография \href{http://mi.mathnet.ru/intm17 }{{\it Олейник~О.~А., Радкевич~Е.~В.} ``Уравнения второго порядка с неотрицательной характеристической формой''/ Итоги науки. Сер. Математика. Мат. анал. 1969. М.: ВИНИТИ,  1971. 252~c.} и~ее переиздания на других языках на многие годы послужили  источником и~вдохновителем исследования вырождающихся уравнений. В расширенное переиздание этой монографии,\footnote{\href{https://elibrary.ru/item.asp?id=19461734}{{\it Олейник~О.~А., Радкевич~Е.~В.} ``Уравнения с неотрицательной характеристической формой''. М.: Моск. ун-т, 2010. 359~с.}}\!\!~связанное с~юбилеем О.~А.~Олейник, были включены работы по гипоэллиптичности и методу введения параметра для исследования эволюционных уравнений. Это издание вызвало новый интерес к проблематике вырождающихся уравнений, что отмечается рядом ссылок на него.



Монография \href{https://elibrary.ru/item.asp?id=19453018}{{\it Радкевич~Е.~В.} ``Математические вопросы неравновесных процессов''/ Белая серия в математике и физике. Т.~4. Новосибирск: Изд. Т.~Рожковская, 2007. 300~с.} посвящена классической проблеме замыкания цепочки моментных уравнений кинетического уравнения Больцмана и связанной с ней проблемы обоснования  приближения Навье"--~Стокса. Результаты этой монографии (и исследования, выполненные совместно с~\rupid{8919}{Л.~Р.~Волевичем}, \rupid{28276}{П.~А.~Захарченко}, \rupid{52554}{И.~В.~Загребаевым} и \rupid{102571}{В.~В.~Палиным}) позволили получить условия разрешимости матричных уравнений, обобщения условия Гурвица корректности гиперболических уравнений и~широкое обобщение метода Ньютона для получения условий разрешимости задач со свободной границей.

Книги за авторством Е.~В.~Радкевича переведены и изданы на нескольких языках.







В цикле работ [\citen{III.1,  III.2, III.3, III.6, III.7, III.8}],~посвященных исследованию так называемых дискретных кинетических уравнений (Карлемана, Годунова"--~Султангазина, Браудера), получены условия существования глобального решения задачи Коши для начальных данных  (периодических и~локализованных) возмущений положительных состояний равновесия. 



Получены условия локального равновесия для решения задачи Коши с~локализованными начальными данными и~ограниченной энергией. 

Доказана асимптотическая устойчивость положительных состояний равновесия (с~экспоненциальной стабилизацией).

Для возмущений неустойчивых стационарных решений смешанной задачи доказано возникновение хаотической динамики.



В цикле работ~\cite{Add, IV.2, IV.4, IV.5, IV.7, IV.8, IV.9, IV.10,  IV.12, IV.13, IV.14, IV.15, IV.16, IV.17, IV.18} исследуется класс критических процессов, объединенных общей гипотезой, что они есть неравновесные фазовые переходы (ламинарно"=турбулентный переход, неустойчивость Рэлея"--~Бенара, кристаллизация бинарных сплавов, разрушение конструкционных материалов, неустойчивость Марангони). 



Для ламинарно"=турбулентного перехода, неустойчивости Рэлея"--~Бенара и кристаллизации бинарных сплавов построены модели реконструкции начальной стадии неустойчивости как неравновесного перехода, механизмом которого является диффузионное расслоение.



Показано, что свободная энергия Гиббса отклонения от однородного состояния (относительно рассматриваемой неустойчивости) есть аналог потенциалов Гинзбурга"--~Ландау. 

Проведены численные эксперименты самовозбуждения однородного состояния управлением краевым условием (возрастанием скорости или температуры на входе). 

Установлена нелокальность возмущения, что указывает на невозможность применения в этом случае классической теории возмущения. 



В \cite{Add} установлено, что распространение теории неравновесных фазовых переходов на задачу Бенара"--~Марангони проблематично. 

Первоначальный термодинамический анализ задачи Бенара"--~Марангони установил невозможность построения аналога потенциала Гинзбурга"--~Ландау в~этом случае и~тем самым применения предложенной схемы построения реконструкции для начальной стадии неустойчивости Марангони. 



В работах \cite{IV.13, IV.14, IV.15, IV.12, IV.9, IV.7} используется новый подход к~описанию турбулентности, который имеет важное прикладное значение в~связи с~проблемами гиперзукового движения. 

Были исследованы прикладные задачи, связанные с~проблемой угасания гиперзвукового реактивного двигателя, и~создана математическая модель процесса, в~котором возникают ветвящиеся ударные волны, аналогичные наблюдаемым в~эксперименте.



Е.~В.~Радкевич "--- прекрасный лектор и~опытный педагог. 

На протяжении многих лет Евгений Владимирович читает лекции и~ведет семинары по уравнениям с~частными производными на механико"=математическом факультете и~на факультете фундаментальной физико"=химической инженерии МГУ. 

Под руководством Е.~В.~Радкевича подготовлено и~защищено 6 диссертаций на соискание ученой степени кандидата наук и~одна "--- на соискание степени доктора наук.



Е.~В.~Радкевич является соруководителем популярного активно работающего семинара кафедры дифференциальных уравнений по асимптотическим методам математической физики и~членом диссертационного совета при механико"=математическом факультете МГУ. 

Является постоянным участником множества научных конференций, сотрудничает с ведущими мировыми научными центрами (в~т.ч. \ruaffid{1431}{Институт Вейерштрасса, Германия}). 

Кроме этого, Е.~В.~Радкевич является одним из авторов выдержавшего несколько переизданий задачника по уравнениям с~частными производными,\footnote{\href{https://istina.msu.ru/publications/book/1826310/}{{\it Вентцель~Т.~Д., Горицкий~А.~Ю., Капустина~Т.~О., Кондратьев~В.~А., Радкевич~Е.~В., Розанова~О.~С., Чечкин~Г.~А., Шамаев~А.~С., Шапошникова~Т.~А}. ``Сборник задач по уравнениям с частными производными''. М.: БИНОМ. Лаборатория знаний, 2005. 158~c.}}\!\!~а также одним из авторов учебника по методам математической физики для факультета фундаментальной физико"=химической инженерии МГУ, переизданного в 2017 году.\footnote{

\href{https://elibrary.ru/item.asp?id=30568002}{{\it Палин~В.~В., Радкевич~Е.~В.} ``Методы математической физики. Лекционный курс'': учебное пособие для академического бакалавриата. 2-е изд., испр. и~доп. М.: Изд-во Юрайт, 2017.  222~с.} }



Нельзя не упомянуть и другие интересы Е.~В.~Радкевича, которые он удивительным образом смог совместить со своей научной карьерой. Еще во время учебы он принимал активное участие в работе студенческого театра МГУ под руководством \href{https://goo.gl/YhM9Eg}{П.~П.~Васильева}, затем "--- \href{https://goo.gl/nB1LPU}{П.~Н.~Фоменко}. 

Не прекращая научную и преподавательскую деятельность, в~1969 году поступил, а~в~1976 году окончил \href{https://goo.gl/DfcJNR}{ГИТИС} (специальность <<Режиссура драмы>>, курс \href{https://goo.gl/HFRLHQ}{Андрея Гончарова}, однокурсниками на актерском отделении были Светлана Акимова, \href{https://goo.gl/a8g3Xm}{Александр Соловьев}, \href{https://goo.gl/VRGfi6}{Игорь Костолевский}, \href{https://goo.gl/oQ89tw}{Александр Фатюшин}). Дипломный спектакль <<Рассказ от первого лица>> поставил в \href{https://goo.gl/mvSV2b}{Театре на Малой Бронной} у \href{https://goo.gl/rTdrBN}{Анатолия Эфроса}, позже был режиссером спектаклей в \href{https://goo.gl/uFGAzt}{ТЮЗе}, \href{https://goo.gl/6WLbvz}{Московском Новом драматическом театре}.



Столь разносторонний талант отмечен многочисленными наградами:

\begin{itemize}

\item[--] Премия Московского комсомола; 

\item[--] Лауреат фестиваля <<Московская театральная весна-83>> за спектакль <<Новичок>> (ТЮЗ); 

\item[--] Лауреат Третьего фестиваля молодых кинематографистов Центральной киностудии им.~М.~Горького, 1984 за фильм <<Парамон и Аполлинария>> (короткометражный).

\end{itemize}





Перечислить все научные работы Е.~В.~Радкевича в этой короткой заметке нет возможности. Здесь мы приводим лишь список его публикаций за последние пять лет, который дает представление об уровне публикаций Е.~В.~Радкевича.





\newpage



\PersonRadRef



\begin{thebibliography}{99}



\Bibitem{Add}

\by Radkevich~E.~V., Lukashev~E.~A., Sidorov~M.~I., Vasil'eva~O.~A.

\paper Methods of nonlinear dynamics of nonequilibrium processes in fracture mechanics 

\jour Eurasian Journal of Mathematical and Computer Applications

\vol 6

\issue 2 

\yr 2018 

\toappear





\Bibitem{IV.18}

\paper On the Rayleigh–Bénard instability as the nonequilibrium phase transition

\jour AIP Conf. Proc. 

\vol 1910

\issue 1

\papernumber 020017 

\yr 2017

\crossref{10.1063/1.5013954}

\by Lukashev~E.~A., Radkevich~E.~V.,  Yakovlev~N.~N.,  Vasil'eva~O.~A.





%Radkevich E.V., Vasil’eva O.A. GENERATION OF CHAOTIC DYNAMICS AND LOCAL EQUILIBRIUM FOR THE CARLEMAN EQUATION в журнале Journal of Mathematical Sciences, издательство Plenum Publishers (United States), 2017г, том 224, № 5, с. 763-795 

\RBibitem{III.1}

\by Радкевич~Е.~В., Васильева~О.~А.

\paper Возникновение хаотической динамики и локальное равновесие для уравнения Карлемана  

\pages 143--170

\jour Пробл. мат. анал.

\vol 87

\yr 2017

\transl

\by Radkevich~E.~V., Vasil’eva~O.~A.

\paper Generation of chaotic dynamics and local equilibrium for the Carleman equation

\jour J. Math. Sci.

\yr 2017

\vol 224

\issue 5

\pages 763--795

\crossref{10.1007/s10958-017-3449-6}



\Bibitem{IV.16}

\paper  	\href{http://mi.mathnet.ru/ejmca18 }{Study of the Rayleigh-Benard instability by methods of the theory of nonequilibrium phase transitions in the Cahn--Hillard form}

\by Radkevich~E.~V.,  Lukashev~E.~A., Yakovlev~N.~N., Vasil'eva~O.~A.

\jour Eurasian Journal of Mathematical and Computer Applications

\vol 5

\issue 2 

\yr 2017

\pages 36–65





\RBibitem{IV.2}

\by Лукашев~Е.~А., Радкевич~Е.~В., Яковлев~Н.~Н., Васильева~О.~А.

\paper Введение в обобщенную теорию неравновесных фазовых переходов Кана—Хилларда

(термодинамический анализ задач механики сплошной среды)

\jour Вестн. Сам. гос. техн. ун-та. Сер. Физ.-мат. науки

\yr 2017

\vol 21

\issue 3

\pages 437--472

\mathnet{http://mi.mathnet.ru/vsgtu1554}

\crossref{10.14498/vsgtu1554}

\elib{http://elibrary.ru/item.asp?id=32248390}





%Палин В.В., Радкевич Е.В. О поведении стабилизирующихся решений для уравнения Риккати в журнале Сборник "Труды семинара им. И.Г. Петровского", 2017г, том 31, с. 111-134 

\RBibitem{PalRad16}

\by Палин~В.~В., Радкевич~Е.~В.

\paper \href{{http://mi.mathnet.ru/tsp92}}{О~поведении стабилизирующихся решений для уравнения Риккати}

\serial Тр. сем. им. И.~Г.~Петровского

\yr 2016

\vol 31

\pages 110--133

\publ Изд-во Моск. ун-та

\publaddr М.

\mathnet{http://mi.mathnet.ru/tsp92}



%Радкевич Е.В., Лукашев Е.А., Яковлев Н.Н., Васильева О.А. О природе конвективной неустойчивости Рэлея-Бенара в журнале Доклады Академии наук, издательство Наука (М.), 2017г, том 475, № 6, с. 618-623 

\RBibitem{IV.17}

\by Радкевич~Е.~В., Лукашев~Е.~А., Яковлев~Н.~Н., Васильева~О.~А.

\paper \href{https://elibrary.ru/item.asp?id=29766739}{О природе конвективной неустойчивости Рэлея--Бенара}

\jour Докл. Акад. наук

\yr 2017

\vol 475

\issue 6

\pages 618--623

\transl

\paper On the nature of the Rayleigh–Bénard convective instability 	

\by \linebreak Radkevich~E.~V., Lukashev~E.~A., Yakovlev~N.~N., Vasil'eva~O.~A.

\yr 	2017

\jour	Dokl. Math.

\vol 96

\issue 1

\pages 393--398 

\crossref{10.1134/S1064562417040317}





%Lukashev E.A., Yakovlev N.N., Radkevich E.V., Vasil'yeva O.A. On the theory of nonequilibrium phase transitions on the laminar-turbulent transition в журнале Наноструктуры. Математическая физика и моделирование, 2016г., том 14, № 1, с. 5-40

\RBibitem{IV.13}

\by Лукашев~Е.~А., Яковлев~Н.~Н., Радкевич~Е.~В., Васильева~О.~А.

\jour Наноструктуры. Математическая физика и моделирование

\paper \href{https://elibrary.ru/item.asp?id=26711248}{О распространении теории неравновесных фазовых переходов на ламинарно-турбулентный переход}

\yr 2016

\vol 14

\issue 1 

\pages 5--40







%ВАСИЛЬЕВА О.А., ДУХНОВСКИЙ С.А., РАДКЕВИЧ Е.В. О ПРИРОДЕ ЛОКАЛЬНОГО РАВНОВЕСИЯ УРАВНЕНИЙ КАРЛЕМАНА И ГОДУНОВА—СУЛТАНГАЗИНА в журнале Современная математика. Фундаментальные направления, 2016г., том 60, с. 23-81 

\RBibitem{1}

\by Васильева~О.~А., Духновский~С.~А., Радкевич~Е.~В.

\paper \href{http://mi.mathnet.ru/cmfd295}{О природе локального равновесия уравнений Карлемана и Годунова--Султангазина}

\inbook Труды Седьмой Международной конференции по дифференциальным и функционально-дифференциальным уравнениям \rm (Москва, 22--29 августа, 2014). Часть~3

\serial СМФН

\yr 2016

\vol 60

\pages 23--81

\publ РУДН

\publaddr М.

\mathnet{http://mi.mathnet.ru/cmfd295}



%Лукашев Е.А., Яковлев Н.Н., Радкевич Е.В., Васильева О.А. О ПРОБЛЕМАХ ЛАМИНАРНО-ТУРБУЛЕНТНОГО в журнале Доклады Академии наук, издательство Наука (М.), 2016г., том 471, № 3, с. 1-5 

\RBibitem{IV.14}

\by Лукашев~Е.~А., Яковлев~Н.~Н., Радкевич~Е.~В., Васильева~О.~А.

\paper О проблемах ламинарно-турбулентного перехода

\jour Докл. Акад. наук

\yr 2016

\vol 471

\issue 3

\pages 270--274

\crossref{10.7868/S0869565216330045}

\transl

\paper On problems of the laminar–turbulent transition

\by Lukashev~E.~A., Yakovlev~N.~N., Radkevich~E.~V., Vasil'yeva~O.~A.

\jour Dokl. Math.

\yr  2016

\vol 94

\issue 3

\pages 649--653

\crossref{10.1134/S1064562416060119}





%Лукашев Е.А., Радкевич Е.В., Яковлев Н.Н., Васильева О.А. О реконструкции начальной стадии турбулентно-диффузионного горения в журнале Сибирский журнал чистой и прикладной математики, 2016г., № 2, с. 50-67 

\RBibitem{IV.15}

\by Лукашев~Е.~А., Радкевич~Е.~В., Яковлев~Н.~Н., Васильева~О.~А.

\paper О реконструкции начальной стадии турбулентно-диффузионного горения

\jour Сиб. журн. чист. и прикл. матем.

\yr 2016

\vol 16

\issue 2

\pages 50--67

\mathnet{http://mi.mathnet.ru/vngu402}

\crossref{10.17377/PAM.2016.16.205}





\RBibitem{III.2}

\paper \href{https://elibrary.ru/item.asp?id=28111054}{О локальном равновесии уравнения Карлемана}

\by Васильева~О.~А., Духновский~С.~А., Радкевич~Е.~В.

\jour Пробл. мат. анал.

\yr 2015

\vol 78

\pages 165--190

\transl

\paper Local equilibrium of the Carleman equation

\by Radkevich~E.~V., Vasil'eva~O.~A., Dukhnovskii~S.~A.

\jour J. Math. Sci. 

\yr 2015

\vol 207

\issue 2

\pages 296--323

\crossref{10.1007/s10958-015-2373-x}





%Radkevich E.V., Palin V.V. On the Riemann–Hugoniot Catastrophe в журнале Russian Journal of Mathematical Physics, издательство Maik Nauka/Interperiodica Publishing (Russian Federation),2015г., том 22, № 2, с. 10-23 

\Bibitem{1} 

\paper On the Riemann--Hugoniot catastrophe 	

\by Palin~V.~V., Radkevich~E.~V.

\yr 	2015

\jour Russ. J. Math. Phys. 

\yr 2015

\vol 22

\issue 2

\crossref{10.1134/S1061920815020090}

\pages 227--236 



%Palin V.V., Radkevich E.V. On the nature of bifurcations of solutions of the Riemann problem for the truncated Euler system в журнале Differential Equations, издательство Maik Nauka/Interperiodica Publishing (Russian Federation), 2015г., том 51, № 6, с. 755-766 DOI 

\RBibitem{1}

\paper О природе бифуркаций решений задачи Римана для усеченной системы Эйлера

\by Палин~В.~В., Радкевич~Е.~В.

\jour Диффер. урав.

\yr  2015

\vol 51

\issue 6

\pages 743--754

\crossref{10.1134/S0374064115060060}

\transl

\paper On the nature of bifurcations of solutions of the Riemann problem for the truncated Euler system

\by Palin~V.~V., Radkevich~E.~V.

\jour Differ. Equat.

\yr  2015

\vol 51

\issue 6

\pages 755--766

\crossref{10.1134/S0012266115060063}





\Bibitem{1}

\by Radkevich~E.~V.

\paper The Bloch principle and $L_2 (R)$ local equilibrium of the Carleman equation

\jour AIP Conf. Proc.

\vol 1690

\issue 

\papernumber 040006 

\yr 2015

\crossref{10.1063/1.4936713}







%Radkevich E.V.THE BLOCH PRINCIPLE FOR L 2 (R) STABILIZATIONOF SOLUTIONS TO THE CAUCHY PROBLEM FORTHE CARLEMAN EQUATION в журнале Journal of Mathematical Sciences, издательство Plenum Publishers (United States),2015г., том 210, № 5 

\Bibitem{1}

\by Radkevich~E.~V.

\paper The Bloch principle for $L_2 (R)$ stabilizationof solutions to the Cauchy problem for the Carleman equation

\jour J. Math. Sci.

\yr 2015

\vol 210

\issue 5

\crossref{10.1007/s10958-015-2586-z}

\pages 677–735 



%Яковлев Н.Н., Лукашёв Е.А., Радкевич Е.В., Палин В.В.  О парадигме внутренней турбулентности в журнале Вестник СамГТУ, серия «Физико- математические науки», 2015г., том 19, № 1, с. 155-185 DOI 

\RBibitem{IV.12}

\by Яковлев~Н.~Н., Лукашев~Е.~А., Радкевич~Е.~В., Палин~В.~В. 

\paper О парадигме внутренней турбулентности

\jour Вестн. Сам. гос. техн. ун-та. Сер. Физ.-мат. науки

\yr 2015

\vol 19

\issue 1

\pages 155--185

\crossref{10.14498/vsgtu1418}





%Палин В.В., Радкевич Е.В., Яковлев Н.Н., Лукашёв Е.А. О невязких решениях многокомпонентной системы Эйлера в журнале Современная математика. Фундаментальные направления, 2014г., том 53, с. 133-154 

\RBibitem{IV.8}

\by Палин~В.~В., Радкевич~Е.~В., Яковлев~Н.~Н., Лукашев~Е.~А.

\paper \href{http://mi.mathnet.ru/cmfd263}{О невязких решениях многокомпонентной системы Эйлера}

\inbook Труды Крымской осенней математической школы-симпозиума

\serial СМФН

\yr 2014

\vol 53

\pages 133--154

\publ РУДН

\publaddr М.

\mathnet{http://mi.mathnet.ru/cmfd263}

\transl

\paper On nonviscous solutions of a multicomponent Euler system

\by  Palin~V.~V.,  Radkevich~E.~V.,  Yakovlev~N.~N., Lukashev~E.~A.

\jour J. Math. Sci.

\yr 2016

\vol 218

\issue 4

\pages 503--525

\crossref{10.1007/s10958-016-3040-6}





%Асташова И.В., Боровских А.В., Быков В.В., Ветохин А.Н., Горицкий А.Ю., Изобов Н.А., Ильяшенко Ю.С., Капустина Т.О., Козлов В.В., Коньков А.А., Матросов И.В., Палин В.В., Розов Н.Х., Романов М.С., Сергеев И.Н., Радкевич Е.В., Розанова О.С., Филимонова И.В., Филиновский А.В., Чечкин Г.А., Шамаев А.С., Шапошникова Т.А. Научное наследие Владимира Михайловича Миллионщикова в журнале Труды семинара имени И.Г.Петровского, 2014г., № 30, с. 5-41 

\RBibitem{1}

\by Асташова~И.~В., Боровских~А.~В., Быков~В.~В., Ветохин~А.~Н., Горицкий~А.~Ю., Изобов~Н.~А., Ильяшенко~Ю.~С., Капустина~Т.~О., Козлов~В.~В., Коньков~А.~А., Матросов~И.~В., Палин~В.~В., Розов~Н.~Х., Романов~М.~С., Сергеев~И.~Н., Радкевич~Е.~В., Розанова~О.~С., Филимонова~И.~В., Филиновский~А.~В., Чечкин~Г.~А., Шамаев~А.~С., Шапошникова~Т.~А. 

\serial Тр. сем. им. И.~Г.~Петровского

\paper \href{http://mi.mathnet.ru/tsp68}{Научное наследие Владимира Михайловича Миллионщикова}

\yr 2014

\vol 30

\pages 5--41

\publ Изд-во Моск. ун-та

\publaddr М.

\mathnet{http://mi.mathnet.ru/tsp68}

\transl

\jour J. Math. Sci. 

\yr 2015

\vol 210

\issue 2

\pages 115--134

\crossref{10.1007/s10958-015-2551-x}



%Радкевич Е.В. НЕВЯЗКИЕ РЕШЕНИЯ МНОГОКОМПОНЕНТНОЙ СИСТЕМЫ ЭЙЛЕРА в журнале Доклады Академии наук, издательство Наука (М.), 2015г.,  том 455, № 4, с. 1-5 

\RBibitem{1}

\by Радкевич~Е.~В.

\paper Невязкие решения многокомпонентной системы Эйлера

\jour Докл. Акад. наук

\yr 2014

\vol 455

\issue 4

\pages 389--393

\crossref{10.7868/S0869565214100065}

\transl

\by Radkevich~E.~V.

\paper Inviscid solutions of the multicomponent Euler system

\jour Dokl. Math.

\yr 2014

\vol 89

\issue 2

\pages 197--201

\crossref{10.1134/S1064562414020227}





%Вентцель Т.Д., Горицкий А.Ю., Капустина Т.О., Кондратьев В.А., Радкевич Е.В., Розанова О.С., Чечкин Г.А., Шамаев А.С., Шапошникова Т.А. Сборник задач по уравнениям с частными производными, место издания БИНОМ. Лаборатория знаний Москва, 2014, ISBN 978-5-94774-721-8 , 158 с.

%\RBibitem{1}

%\by Вентцель~Т.~Д., Горицкий~А.~Ю., Капустина~Т.~О., Кондратьев~В.~А., Радкевич~Е.~В., Розанова~О.~С., Чечкин~Г.~А., Шамаев~А.~С., Шапошникова~Т.~А.

%\book Сборник задач по уравнениям с частными производными

%\publaddr М.

%\publ БИНОМ. Лаборатория знаний 

%\yr 2014

%\totalpages 158

%\transl

%\book A boundary-value problem for a first-order hyperbolic system in a two-dimensional domain

%\by Zhura~N.~A.,  Soldatov~A.~P.

%\publaddr Izv. Math.

%\publ 

%\yr 2014

%\totalpages 158











%Radkevich E.V., Palin V.V. On the Possibility of the Cahn-Hilliard Approach Extension to the Solution of Gas Dynamics Problems (Inner Turbulence) в сборнике 40th International Conference Applications of Mathematics in Engineering and Economics (AMEE'14), серия AIP Conference Proceedings 1631, с. 197-209 DOI 

\Bibitem{IV.7}

\by Radkevich~E.~V., Palin~V.~V.

\paper On the possibility of the Cahn--Hilliard approach extension to the solution of gas dynamics problems (inner turbulence)

\jour AIP Conf. Proc.

\vol 1631

\issue 1

\papernumber 197 

\yr 2014

\crossref{10.1063/1.4902477}



\RBibitem{1}

\by Палин~В.~В., Радкевич~Е.~В.

\paper О бифуркации критических волн разрежения

\jour Пробл. мат. анал.

\yr 2014

\vol 76

\pages 121--132

\transl

\paper  Bifurcations of critical rarefaction waves

\yr 2014 

\jour  J. Math. Sci.

\vol 202 

\issue 2 

\pages 245--258

\by  Palin~V.~V.,  Radkevich~E.~V.

\crossref{10.1007/s10958-014-2044-3}











%Лукашев Е.А., Радкевич Е.В., Яковлев Н.Н.О визуализации начальной стадии кристаллизации бинарных сплавов в журнале Наноструктуры. Математическая физика и моделирование, 2014г., том 11, № 2, с. 5-36 

\RBibitem{IV.4}

\by Лукашев~Е.~А., Радкевич~Е.~В., Яковлев~Н.~Н.

\paper О визуализации начальной стадии кристаллизации бинарных сплавов

\jour Наноструктуры. Математическая физика и моделирование

\yr 2014

\vol 11

\issue 2

\pages 5--36













%Яковлев Н.Н., Лукашев Е.А., Радкевич Е.В. О реконструкции начальной стадии внутренней турбулентности в журнале Наноструктуры. Математическая физика и моделирование, 2014г., том 11, № 1, с. 73-99 

\RBibitem{IV.9}

\by Яковлев~Н.~Н., Лукашев~Е.~А., Радкевич~Е.~В.

\paper \href{https://elibrary.ru/item.asp?id=24258345}{О реконструкции начальной стадии внутренней турбулентности}

\jour Наноструктуры. Математическая физика и моделирование

\yr 2014

\vol 11

\issue 1

\pages 73--100





\RBibitem{1}

\by Радкевич~Е.~В.

\paper О природе биффуркаций однофронтовых решений усеченной системы Эйлера

\jour Пробл. мат. анал.

\yr 2013

\vol 73

\pages 125–139

\transl

\paper On the nature of bifurcations of one-front solutions of the truncated Euler system

\yr 2014 

\jour J. Math. Sci.

\yr 2014

\vol 196 

\issue 3 

\pages 388--404

\by Radkevich~E.~V. 

\crossref{10.1007/s10958-014-1664-y}





\RBibitem{IV.5}

\paper  О неклассической регуляризации многокомпонентной системы Эйлера

\by Лукашев~Е.~А., Палин~В.~В., Радкевич~Е.~В., Яковлев~Н.~Н.

\jour Пробл. мат. анал.

\yr 2013

\vol 73

\pages 67–86

\transl

\jour J. Math. Sci.

\yr 2014

\crossref{10.1007/s10958-014-1661-1}

\vol 96

\issue 3

\pages 322–345 

\paper Nonclassical regularization of the multicomponent Euler system

\by \linebreak Lukashev~E.~A., Palin~V.~V., Radkevich~E.~V., Yakovlev~N.~N.





\RBibitem{III.8}

\by Е.~В.~Радкевич

\paper \href{http://mi.mathnet.ru/cmfd225}{О поведении на больших временах решений задачи Коши для двумерного дискретного кинетического уравнения}

\inbook Труды Шестой Международной конференции по дифференциальным и функционально-дифференциальным уравнениям \rm (Москва, 14--21 августа, 2011). Часть~3

\serial СМФН

\yr 2013

\vol 47

\pages 108--139

\publ РУДН

\publaddr М.

\mathnet{http://mi.mathnet.ru/cmfd225}

\transl

\by Radkevich~E.~V.

\paper On the large-time behavior of solutions to the Cauchy problem for a 2-dimensional discrete kinetic equation

\jour J. Math. Sci.

\yr 2014

\vol 202

\issue 5

\pages 735--768

\crossref{10.1007/s10958-014-2074-x}

\scopus{http://www.scopus.com/record/display.url?origin=inward&eid=2-s2.0-84921936201}



\Bibitem{III.7}

\by Radkevich~E.~V.

\crossref{10.1007/s10958-013-1213-0}

\jour J. Math. Sci.

\yr 2013

\vol  189

\issue 4

\pages 659–698 

\paper On nonexistence of dissipative estimates for discrete kinetic equations





\RBibitem{Rad13}

\by Радкевич~Е.~В.

\paper К~проблеме несуществования диссипативной оценки для~дискретных кинетических уравнений

\jour Вестн. Сам. гос. техн. ун-та. Сер. Физ.-мат. науки

\yr 2013

\issue 1(30)

\pages 106--143

\mathnet{http://mi.mathnet.ru/vsgtu1140}

\crossref{10.14498/vsgtu1140}





\RBibitem{III.3}

\by Загребаев~И.~В., Радкевич~Е.~В.

\paper \href{http://mi.mathnet.ru/cmfd174}{О промежуточных аттракторах}

\inbook Уравнения в частных производных

\serial СМФН

\yr 2011

\vol 39

\pages 79--101

\publ РУДН

\publaddr М.

\mathnet{http://mi.mathnet.ru/cmfd174}

\mathscinet{http://www.ams.org/mathscinet-getitem?mr=2830678}

\transl

\paper On intermediate attractors

\by Zagrebaev~I.~V., Radkevich~E.~V. 

\jour J.~Math. Sci.

\yr 2013

\vol 190

\issue 1

\pages 80--103

\crossref{10.1007/s10958-013-1247-3}

\scopus{http://www.scopus.com/record/display.url?origin=inward&eid=2-s2.0-84874949616}



%Палин В.В., Радкевич Е.В. Методы математической физики (уравнения с частными производными второго порядка)Лекционный курс: методы математической физики. Часть II, место издания Издательство попечительского совета мехмата МГУ Москва, 2012, 104 с. 

%\RBibitem{1}

%\by Палин~В.~В., Радкевич~Е.~В.

%\book Методы математической физики (уравнения с частными производными второго порядка)/ Лекционный курс: методы математической физики. Часть II

%\publaddr М.

%\publ Издательство попечительского совета мехмата МГУ

%\yr 2012

%\totalpages 104







%Палин В.В., Радкевич Е.В. Методы математической физики (уравнения с частными производными первого порядка и коротковолновая асимптотика) Лекционный курс: методы математической физики. Часть I, место издания Издательство попечительского совета мехмата МГУ Москва, 2012, 120 с. 

%\RBibitem{1}

%\by Палин~В.~В., Радкевич~Е.~В.

%\book Методы математической физики (уравнения с частными производными второго порядка)/ Лекционный курс: методы математической физики. Часть I

%\publaddr М.

%\publ Издательство попечительского совета мехмата МГУ

%\yr 2012

%\totalpages 120





%Yakovlev N.N., Lukashev E.A., Radkevich E.V. Investigation of Directed Crystallization by the Method of Mathematical Reconstruction в журнале Doklady Physics, издательство Maik Nauka/Interperiodica Publishing (Russian Federation), 2012г.,  том 57, № 8, с. 297-300 DOI 

\RBibitem{IV.10}

\paper \href{https://elibrary.ru/item.asp?id=17858912}{Исследование процесса направленной кристаллизации методом математической реконструкции}

\by Яковлев~Н.~Н., Лукашев~Е.~А., Радкевич~Е.~В.

\jour Докл. Акад. наук

\yr 2012

\vol 445

\issue 4

\pages 398--401

\transl

\by Yakovlev~N.~N.,  Lukashev~E.~A., Radkevich~E.~V.

\paper Investigation of directed crystallization by the method of mathematical reconstruction

\jour Dokl. Phys.

\yr 2012

\vol 57

\issue 8

\pages 297--300

\crossref{10.1134/S1028335812080046}







%Radkevich E.V. On discrete kinetic equations в журнале Doklady Mathematics, издательство Maik Nauka/Interperiodica Publishing (Russian Federation), 2012г., том 86, № 3, с. 809-813 DOI 

\RBibitem{III.6}

\paper \href{https://elibrary.ru/item.asp?id=18447865}{О дискретных кинетических уравнениях}

\by Радкевич~Е.~В.

\jour Докл. Акад. наук

\yr 2012

\vol 447

\issue 4

\pages 369--373

\transl

\by Radkevich~E.~V.

\paper On discrete kinetic equations

\jour Dokl. Math.

\yr 2012

\vol 86

\issue 3

\pages 809--813

\crossref{10.1134/S1064562412060233}







%Radkevich E.V. THE EXISTENCE OF GLOBAL SOLUTIONS TO THE CAUCHY PROBLEM FOR DISCRETE KINETIC EQUATIONS(NONPERIODIC CASE в журнале Journal of Mathematical Sciences, издательство Plenum Publishers (United States), 2012г., том 184, № 4, с. 524-556 

\Bibitem{1}

\by Radkevich~E.~V.

\paper The existence of global solutions to the Cauchy problem for discrete kinetic equations (nonperiodic case)

\jour J.~Math. Sci.

\yr 2012

\vol 184

\issue 4

\pages 524--556

\crossref{10.1007/s10958-012-0879-z}





\Bibitem{1}

\paper  The existence of global solutions to the Cauchy problem for discrete kinetic equations. II

\by Radkevich~E.~V.

\jour J.~Math. Sci.

\yr 2012

\vol 181

\issue 5

\pages 701--750

\crossref{10.1007/s10958-012-0711-9}





\Bibitem{1}

\paper  The existence of global solutions to the Cauchy problem for discrete kinetic equations 

\by Radkevich~E.~V.

\jour J.~Math. Sci.

\yr 2012

\vol 181

\issue 2

\pages 232--280

\crossref{10.1007/s10958-012-0683-9}2





\RBibitem{1}

\paper   \href{https://elibrary.ru/item.asp?id=17951625}{Об осцилляциях, порождаемых оператором взаимодействия в дискретных кинетических уравнениях}

\by Радкевич~Е.~В.

\jour Научные ведомости БелГУ. Математика. Физика

\yr 2012

\vol 26

\issue 5

\pages 145--2062





\end{thebibliography}





\vfill



Редколлегия журнала <<Вестник Самарского государственного технического университета. Сер. Физико-математические науки>> поздравляет Евгения Владимировича с юбилеем и желает ему здоровья, творческих успехов в научной и педагогической деятельности и счастья.







\vfill







\hfill {\it \rupid{38541}{А.~А.~Андреев}, \rupid{38375}{В.~П.~Радченко}, \rupid{62329}{Е.~А.~Козлова}}





\vfill





\AdditionalContent{Используемые материалы}{В заметке были использованы материалы \href{http://new.math.msu.su/}{научного портала механико-математического факультета МГУ} (\url{http://new.math.msu.su/staff/radkevic.htm}) и Википедии  (\href{https://goo.gl/3xRf6i}{\tt Радкевич, Евгений Владимирович}). Фотография любезно предоставлена юбиляром.} 



\vfill



\vfill







\newpage



\selectlanguage{english}%

\normalsize 

\normalfont

\chead[\fancyplain{}{\selectlanguage{english}\theenauthor}]{\fancyplain{}{\selectlanguage{english} \theentitle}}

\rhead[]{\fancyplain{\EnJournal}{}}

\lhead[\fancyplain{\EnJournal}{}]{}

\cfoot[]{}

\lfoot[\thepage]{}

\rfoot[]{\thepage}

\renewcommand{\plainheadrulewidth}{0.7pt}

\renewcommand{\headrulewidth}{0.7pt}

\pagestyle{fancyplain}

\thispagestyle{plain}

\clear

\hskip-7mm

\begin{minipage}[t]{125mm}

\vskip.5mm

\noindent\themsc

\vskip2mm

\EnTitleInTitlepage

\end{minipage}

\par

\theenauthorsdetails

\begin{minipage}[t]{125mm}

\vskip4mm

{\renewcommand{\bfdefault}{bx}\noindent\small\textbf{\emph{\EnAuthorInTitlepage}}\renewcommand{\bfdefault}{b}}  

\vskip2mm

{\scriptsize \enaffilname}

\end{minipage}

\vskip2mm





\footnotesize



\smallskip



\selectlanguage{english}



One of the most influential scientists in mathematical science,  Evgeniy Vladimirovich Radkevich, has celebrated his 75$^{\text{th}}$ anniversary on the January, 26$^{\text{th}}$. Here we give the biographical background and discuss his well known researches in different areas of science and technology.



Evgeniy Vladimirovich Radkevich began his mathematical education in Moscow State University, Faculty of Mechanics and Mathematics. He has studied under Olga Arsenievna Oleinik, a famous mathematician, one of the best specialists in the theory of partial differential equations, and has become a Candidate of Physical and Mathematical Sciences. E.~V.~Radkevich worked in the leading mathematical institutes---Institute for Problems in Mechanics of the USSR Academy of Sciences, Moscow Institute of Electronics and Mathematics, Moscow Institute of Radio-Engineering, Electronics and Automation, Moscow State University. 



Now he is Professor, Doctor habilitated of Physical and Mathematical Sciences, his interests are in many areas of mathematics (differential equations and boundary value problems, asymptotical methods, phase transition problems and so on) and also he is an experienced teacher and lecturer. 



In the paper we give a list of Evgeniy Vladimirovich publications for the last five years to show the level of his scientific researches. 

Also here we discuss his greatest scientific and technological contributions and concern the other important area of interest of  E.~V.~Radkevich---the theatre arts, and his work as theatre director.





\bigskip



\par\noindent

\theonlineengdate \par\noindent

\rule{125mm}{.7pt}\\



\footnotetext{

\vspace{-1em}\\

\noindent{\normalsize\bf\enArticleType}\smallskip\par\noindent

\OAlogo~\ccby\ The content is published under the terms of the \href{http://creativecommons.org/licenses/by/4.0/}{Creative Commons \mbox{Attribution} 4.0 International} License (\url{http://creativecommons.org/licenses/by/4.0/})\selectlanguage{english}

\vspace{.3em}\\           

\EnCitation}









\selectlanguage{russian}





\newpage



%\end{document}

%

